\documentclass[11pt]{article}
\usepackage[margin=1in]{geometry}
\usepackage{amsmath, amssymb, amsthm}
\usepackage{enumitem}
\usepackage{xcolor}
\usepackage{hyperref}
\hypersetup{colorlinks=true, linkcolor=blue!50!black}

\theoremstyle{definition}
\newtheorem{defn}{Definition}[section]
\newtheorem{thm}[defn]{Theorem}
\newtheorem{lem}[defn]{Lemma}
\newtheorem{ex}[defn]{Example}
\newtheorem{rem}[defn]{Remark}

\newcommand{\N}{\mathbb{N}}
\newcommand{\R}{\mathbb{R}}
\newcommand{\eps}{\varepsilon}
\newcommand{\key}[1]{\textcolor{red!70!black}{\textbf{#1}}}

\title{MATH 3150 --- Notes}
\author{}
\date{}

\begin{document}
\maketitle

% =====================================================================
\section{Subsequences}

\begin{defn}[Subsequence]
Let $(x_n)_{n=1}^\infty$ be a sequence of real numbers, and let
$(n_k)_{k=1}^\infty$ be a sequence of natural numbers that is
\key{strictly increasing}:
\[
  n_1 < n_2 < n_3 < \cdots .
\]
The sequence $(x_{n_k})_{k=1}^\infty = (x_{n_1}, x_{n_2}, x_{n_3}, \dots)$
is called a \textbf{subsequence} of $(x_n)$.
\end{defn}

Equivalently, a subsequence is $x \circ \varphi$, where
$x:\N\to\R$ is the sequence and $\varphi:\N\to\N$, $\varphi(k)=n_k$,
is strictly increasing.

\begin{rem}
\leavevmode
\begin{itemize}[nosep]
  \item The \emph{index} of a subsequence is $k$, not $n$. The $k$-th term
        of the subsequence is $x_{n_k}$.
  \item A subsequence keeps terms \textbf{in their original order} and
        never repeats an index. You may skip terms, but you may not reorder
        them or reuse them.
  \item Every sequence is a subsequence of itself (take $n_k=k$).
\end{itemize}
\end{rem}

\begin{lem}\label{lem:nk}
If $(n_k)$ is a strictly increasing sequence in $\N$, then $n_k \ge k$ for
all $k\in\N$.
\end{lem}

\begin{proof}
Induction on $k$. For $k=1$: $n_1\in\N$, so $n_1\ge 1$. If $n_k\ge k$, then
$n_{k+1}>n_k\ge k$, and since $n_{k+1}$ is an integer, $n_{k+1}\ge k+1$.
\end{proof}

\begin{ex}
Let $x_n = (-1)^n$. Then
\begin{itemize}[nosep]
  \item $n_k = 2k$ gives the subsequence $x_{2k} = 1,1,1,\dots$;
  \item $n_k = 2k-1$ gives the subsequence $x_{2k-1} = -1,-1,-1,\dots$.
\end{itemize}
\end{ex}

\begin{ex}
Let $x_n = \tfrac{1}{n}$. With $n_k = k^2$, the subsequence is
$x_{k^2} = \tfrac{1}{k^2}$, i.e.\ $1, \tfrac14, \tfrac19, \dots$.
\end{ex}

\begin{ex}[Non-examples]
For $x_n=\tfrac1n$:
\begin{itemize}[nosep]
  \item $(\tfrac12, 1, \tfrac14, \tfrac13, \dots)$ is \emph{not} a
        subsequence, because the terms are out of order ($n_1=2>n_2=1$).
  \item $(1, 1, \tfrac12, \tfrac12, \dots)$ is \emph{not} a subsequence,
        because indices repeat ($n_1=n_2=1$).
\end{itemize}
\end{ex}

\begin{thm}
If $x_n \to L$, then every subsequence $(x_{n_k})$ also converges to $L$.
\end{thm}

\begin{proof}
Let $\eps>0$. Since $x_n\to L$, there is $N$ with $|x_n-L|<\eps$ for all
$n\ge N$. If $k\ge N$, then $n_k\ge k\ge N$ by Lemma~\ref{lem:nk}, so
$|x_{n_k}-L|<\eps$.
\end{proof}

\noindent\key{Consequence.} If a sequence has two subsequences with
different limits, it diverges. For example, $(-1)^n$ diverges because
$x_{2k}\to 1$ and $x_{2k-1}\to -1$.


\subsection{Notes}
We need to use definitnos to set up our proofs.

\end{document}
